\chapter[termodeaberturadoprojeto]{Termo de Abertura do Projeto}

%\textcolor{red}{Termo de abertura do projeto / Project Charter. Um documento publicado pelo iniciador ou patrocinador do projeto que autoriza formalmente a existência de um projeto e fornece ao gerente do projeto a autoridade para aplicar os recursos organizacionais nas atividades do projeto.}

\section{Dados do projeto}
\begin{description}
    \item [Nome do Projeto:] CARRINHO DE ENTREGAS AUTÔNOMAS 
    \item [Código:] 1-A
    \item [Patrocinador:] Universidade de Brasília
    % \item [Responsável:] ???
    \item [Gerente:] Guilherme Richard/222006347/richard.evz.99@gmail.com/(61) 98286-2036
\end{description}

\section{Objetivos}
%
%\textcolor{red}{O que a empresa pretende obter com a realização do projeto. Descrever o que se pretende realizar para resolver o problema central ou explorar a oportunidade identificada. Para a correta definição do objetivo siga a regra "SMART":
%\begin{description}

\textcolor{black}{O objetivo específico do projeto é desenvolver e implementar um protótipo de carrinho autônomo. Este veículo será responsável pelo transporte de materiais de até 1 kg entre três departamentos estratégicos: Produção, Qualidade e Engenharia. A operação será gerenciada por uma interface que permitirá aos colaboradores solicitar e rastrear o progresso da entrega em tempo real, otimizando o de trabalho interno.}

\textcolor{black}{Para garantir que o sucesso do projeto seja avaliado de forma objetiva, foram definidos três indicadores-chave de desempenho (KPIs). O sistema deverá alcançar uma redução de 40 no tempo médio de espera por materiais, operar com uma taxa de sucesso de 95 nas missões de entrega sem intervenção humana, e liberar 50 horas mensais de trabalho manual para serem realocadas em atividades de maior valor.}

\textcolor{black}{Este objetivo foi estabelecido e validado através de um consenso entre todas as partes interessadas. Os principais departamentos, incluindo P\&D, Produção, Manutenção, Financeiro e TI, participaram da definição dos requisitos e aprovaram o escopo do projeto, garantindo o alinhamento completo com as metas estratégicas e operacionais.}

\textcolor{black}{A viabilidade do projeto é assegurada por seu escopo focado e pelos recursos disponíveis. A operação em um ambiente interno controlado elimina complexidades, enquanto a utilização de tecnologias consolidadas, como LiDAR e ROS, mitiga os riscos de desenvolvimento. O orçamento e a equipe de engenharia já alocados são adequados para cumprir todas as fases propostas.}

\textcolor{black}{O projeto está estritamente limitado a um cronograma de 4 meses, com data de conclusão final fixada para Dezembro de 2025. Este prazo total engloba todas as fases essenciais do projeto, desde o desenvolvimento inicial e a montagem do protótipo, passando pelos testes em ambiente controlado, até a análise final dos resultados e a entrega da documentação.}
     
   
   
   % \item [\textit{Specific} (específico):] Deve ser redigido de forma clara, concisa e compreensiva;
   % \item [\textit{Measurable} (mensurável):] O objetivo específico deve ser mensurável, ou seja, 
   %possível de ser medido por meio de um ou mais indicadores;
   % \item [\textit{Agreed} (acordado):] Deve ser acordado com as partes interessadas, ou seja, as áreas envolvidas na empresa: P\&D, Produção, Comercial, Marketing, Financeira, Jurídica, Manutenção, ambiental, entre outras;
    %\item [\textit{Realistic} (realista):] Deve estar centrado na realidade, no que é possível de ser feito considerando as premissas e restrições existentes, como: orçamento e tempo;
    %\item [\textit{Time Bound} (Limitado no tempo):] Deve ter um prazo determinado para sua finalização.
%\end{description}
%}

\section{Mercado-alvo}

O carrinho de entregas autônomas é uma solução altamente versátil, focada em entregar eficiência e significativa redução de custos operacionais. Ele beneficia diretamente o Varejo e E-commerce, onde atua assumindo a distribuição de pacotes pequenos e médios para otimizar a frota de veículos humanos, aliviando a parte mais cara da cadeia logística. No Food Service e Plataformas de Delivery, sua aplicação garante maior agilidade e oferece uma modalidade de entrega sem contato, ideal para shoppings ou ambientes urbanos densos. Internamente, em hospitais e indústrias, o carrinho é vital para a logística, transportando produtos críticos, medicamentos ou peças de reposição com segurança. Por fim, em ambientes residenciais como condomínios, ele melhora o serviço ao levar encomendas da portaria diretamente à porta do destinatário, valorizando a gestão do condomínio. Dessa forma, o produto é uma plataforma de automação essencial para qualquer organização que busque economizar no transporte de curta distância, por um preço acessível e aumentando a precisão de seus fluxos de trabalho.

%\textcolor{red}{Pessoas, empresas, instituições etc. que usufruirão dos produtos, serviços e resultados gerados pelo projeto, cujos requisitos (tópico abaixo) devem atender as suas necessidades. Podem ser internas ou externas à organização, mas, merecem destaque especial, pois, o projeto está sendo feito para atendê-los de forma direta ou indireta.}

\section{Requisitos}

Os fatos essenciais para a aquisição do carrinho de entregas autônomas são direcionados na promessa de eficiência transformadora e otimização de custos. O primeiro e mais poderoso motivador é a redução drástica dos custos operacionais, pois o carrinho substitui ou complementa a mão de obra humana em tarefas repetitivas de curta distância, operando 24/7 sem fadiga e reduzindo despesas contínuas. Associada a isso, o aumento da eficiência e produtividade é notável, já que a navegação autônoma permite o planejamento de rotas otimizado e a eliminação de erros comuns. Outro ponto crucial é a maior segurança : equipado com sensores, o carrinho minimiza acidentes por erro humano e assegura que as mercadorias, armazenadas em compartimentos de acesso restrito, cheguem intactas ao destinatário correto. Por fim, a decisão de compra é reforçada pelo posicionamento de inovação e sustentabilidade, já que o veículo é elétrico e permite a otimização do humano ao liberar colaboradores para funções mais estratégicas e de maior valor agregado.



%\textcolor{red}{Listar os fatos que são essenciais para o consumidor adquirir o produto, exemplo: cor, tamanho, material, tempo de vida-útil, dispositivo de segurança \textbf{x}, entre outros.}

\section{Justificativa}

O projeto do carrinho de entregas autônomas é justificado pela necessidade premente de resolver o persistente alto custo. Atualmente, esse segmento é marcado por altos gastos com mão de obra, combustível, seguro e a ineficiência causada por erros humanos, congestionamentos e a dificuldade de realizar entregas individuais de forma econômica. A oportunidade reside em introduzir uma solução de automação elétrica e contínua que opera 24/7 sem fadiga, minimizando o erro, otimizando rotas de curta distância e liberando os recursos humanos. Além disso, o projeto capitaliza a crescente demanda do mercado por rapidez, rastreabilidade e entregas sem contato, alinhando a operação logística das empresas com as expectativas do consumidor moderno . Em suma, o carrinho autônomo é a resposta tecnológica para transformar o maior problema de custos da logística em um diferencial de eficiência e serviço.

%\textcolor{red}{Informar o problema ou a oportunidade (necessidade) que justifica o porquê de o projeto ser realizado. Por exemplo: atende uma demanda específica do consumidor final; supre uma necessidade do mercado comercializador; é um diferencial X para o órgão regulamentador.}

\section{Indicadores}

O mercado consumidor do carrinho de entregas autônomas é determinado por fatores que mostram a necessidade de automação e a de operações logísticas. A demanda é confirmada pelo Volume de Entregas de Curta Distância, que mede o potencial de substituição da mão de obra humana em rotas urbanas e pela Concentração de Centros de Distribuição Urbanos, indicando o número de pontos de origem de entregas. O potencial de uso em ambientes internos é mapeado pelos Usuários em Ambientes Controlados, como em grandes condomínios ou campos universitários. A viabilidade financeira do produto é diretamente influenciada pelo Custo Médio da Mão de Obra Logística, fundamental para calcular o Retorno sobre o Investimento (ROI).  Outros indicadores de demanda incluem o Número de Pedidos Express ou Same-Day Delivery, que exigem a agilidade do carrinho, e o Número de Empresas de Food Service e Varejo com Delivery Próprio, que buscam maior controle e redução de custos com intermediários. Dessa forma, a maturidade do mercado é avaliada pelo Nível de Automação Existente na Intralogística dos potenciais clientes, que aponta a necessidade de maior precisão, e pelo Valor do Mercado de Logística em crescimento, confirmando a tendência favorável à tecnologia.

%\textcolor{red}{Listar até 10 indicadores que determinam o mercado consumidor do produto desenvolvido: exemplo: 1) n° de alunos da FGA que utilizam ônibus às 18:00; 2) n° de usuários do restaurante universitários, 3) número de idosos classificados como público-alvo no DF e no estado de Goiás, 4) n° de empresas de segurança registradas no DF etc.}
% \section{Aprovações}

% \begin{tabular}{ l l }
%   \textbf{Patrocinador:} & \_\_\_\_\_\_\_\_\_\_\_\_\_\_\_\_\_\_\_\_\_\_\_\_\_\_\_\_\_\_\_\_\_\_ \\
%   & \\
%   \textbf{CEO:} & \_\_\_\_\_\_\_\_\_\_\_\_\_\_\_\_\_\_\_\_\_\_\_\_\_\_\_\_\_\_\_\_\_\_ \\
%   & \\
%   \textbf{Gerente:} & \_\_\_\_\_\_\_\_\_\_\_\_\_\_\_\_\_\_\_\_\_\_\_\_\_\_\_\_\_\_\_\_\_\_ \\
% \end{tabular}





